\chapter{Results: The Duel}\label{ch:duel}

The duel network plays two-snake Battlesnake on the standard $11\times11$ board. It
trained for 2.5 hours: 1{,}109 iterations, 36.3 million game turns.

\section{Learning curve}

\begin{center}
\begin{tikzpicture}
\begin{axis}[reportaxis, xmin=0, xmax=37, ymin=0, ymax=100,
  xlabel={self-play game turns (millions)}, ylabel={win rate (\%)},
  legend style={at={(0.5,-0.24)},anchor=north,legend columns=2,/tikz/every even column/.append style={column sep=1.2em}}]
\addplot[color=cmured] table[x=steps, y=flood, col sep=comma] {data/duel_eval.csv};
\addlegendentry{vs flood-fill bot}
\addplot[color=inkblue] table[x=steps, y=greedy, col sep=comma] {data/duel_eval.csv};
\addlegendentry{vs greedy bot}
\addplot[color=rulegrey, dashed, line width=0.9pt, forget plot] coordinates {(0,74) (37,74)};
\node[font=\scriptsize\itshape, color=cmudark, anchor=north west] at (axis cs:14,73) {flood-fill vs greedy: 74\%};
\addplot[color=amber, densely dotted, line width=1pt, forget plot] coordinates {(33.59,0) (33.59,100)};
\node[font=\scriptsize\itshape, color=amber!80!black, anchor=east, align=right] at (axis cs:33.2,40) {selected\\checkpoint};
\end{axis}
\end{tikzpicture}
\captionof{figure}{Win rate of the duel network during training, measured every 25
iterations over 200 games against each bot (deterministic play). The grey line is the
flood-fill bot's own win rate against greedy, for scale.}
\label{fig:duel-curve}
\end{center}

The curve has three phases.
\begin{itemize}
\item \textbf{Survival} (0--4M turns). The network first learns not to die. Against the
  greedy bot it starts at 54\% almost immediately, because greedy plays into dead ends
  that a cautious random walker avoids. Against flood-fill it starts at 11\%.
\item \textbf{Competence} (4--10M). The win rate against flood-fill climbs to about
  80\% by 10.6M turns. From 9M turns on, the network beats flood-fill more often than
  flood-fill beats greedy.
\item \textbf{Refinement} (10--36M). Slow gains to around 90\% against both bots, with
  checkpoint-to-checkpoint noise of a few points, about what 200-game evaluations
  produce.
\end{itemize}

\section{Training dynamics}

\begin{center}
\begin{minipage}{0.49\linewidth}
\begin{tikzpicture}
\begin{axis}[reportaxis, width=\linewidth, height=5cm, xmin=0, xmax=37, ymin=0,
  xlabel={turns (millions)}, ylabel={mean game length},
  title={\small\color{cmudark}Self-play game length}]
\addplot[color=cmured] table[x=steps, y=turns, col sep=comma] {data/duel_train.csv};
\end{axis}
\end{tikzpicture}
\end{minipage}\hfill
\begin{minipage}{0.49\linewidth}
\begin{tikzpicture}
\begin{axis}[reportaxis, width=\linewidth, height=5cm, xmin=0, xmax=37, ymin=0, ymax=1.2,
  xlabel={turns (millions)}, ylabel={policy entropy (nats)},
  title={\small\color{cmudark}Policy entropy}]
\addplot[color=inkblue] table[x=steps, y=entropy, col sep=comma] {data/duel_train.csv};
\addplot[color=rulegrey, dashed, line width=0.9pt] coordinates {(0,1.0986) (37,1.0986)};
\node[font=\scriptsize\itshape, color=cmudark, anchor=north east] at (axis cs:36.5,1.07) {uniform over 3 moves, $\ln 3$};
\end{axis}
\end{tikzpicture}
\end{minipage}
\captionof{figure}{Left: mean length of self-play games, which grows from about 30 to
about 185 turns as both sides learn to survive. Right: entropy of the masked policy,
falling from 0.95 nats (close to uniform over the usual three safe moves) to about 0.3
nats as play becomes decisive.}
\label{fig:duel-dynamics}
\end{center}

Self-play games lengthening sixfold is the clearest sign that both copies of the network
are improving together: neither side finds easy kills against itself any more. Entropy
falls steadily but never collapses. The $0.01$ entropy bonus keeps enough randomness
that exploration continues to the end of the run.

\section{Final evaluation}

The checkpoint with the best score (win rate $+\,\tfrac12$ draw rate against flood-fill)
was iteration 1{,}025, at 33.6M turns. It was then evaluated on 1{,}000 \emph{fresh}
games against each bot, with a seed never used during training or selection.

\begin{center}
\begin{tabular}{@{}lcccccc@{}}
\toprule
\textbf{Opponent} & \textbf{W} & \textbf{D} & \textbf{L} & \textbf{Win rate (95\% CI)} &
  \textbf{Game length} & \textbf{Final length}\\
& & & & & \footnotesize turns & \footnotesize net\ /\ bot\\
\midrule
flood\_fill & 914 & 4 & 82 & \textbf{91.4\%} \ (89.5--93.0) & 99.3 & 13.1 / 12.5\\
greedy & 902 & 64 & 34 & \textbf{90.2\%} \ (88.2--91.9) & 63.1 & \hphantom{1}9.4 / 10.5\\
safe\_random & 994 & 2 & 4 & \textbf{99.4\%} \ (98.7--99.7) & 36.7 & \hphantom{1}8.3 / \hphantom{1}4.0\\
\bottomrule
\end{tabular}
\end{center}

\begin{keybox}[It wins without out-eating its opponent]
Against the greedy bot the network wins 90\% of games while ending them \emph{shorter}
than its opponent, at 9.4 segments against 10.5. It is not winning the race for food. It
wins by position: cutting off space, steering the opponent toward walls and its own
body, and declining head-to-heads it would lose. These are exactly the ideas built into
the flood-fill bot by hand, and that bot also beats greedy. The network found them from a
$\pm1$ signal, and against flood-fill itself it wins 91\% of games.
\end{keybox}

\begin{warnbox}[Draws against greedy]
6.4\% of games against greedy are draws, against 0.4\% for flood-fill. Greedy charges
straight at food; when both snakes are equally long and arrive at the same cell, both
die. The network does not always avoid this; it treats a draw as better than the
risk of losing. A reward that penalised draws more than $0$ would push it further.
\end{warnbox}
